% ============================================
% MEM-24 (ZH)
% 中文翻译 — SPEC-935-R356
% ============================================
\section*{\textbf{MEM-24} --- 哭泣的护士}
\markboth{MEM-24 --- 哭泣的护士}{}

\noindent \textit{小小的善举}

\noindent \textbf{1964 —— 小小的善举}

\vspace{0.5em}\noindent\rule{\textwidth}{0.5pt}\vspace{0.5em}

\noindent 时不时，会有一个新护士出现。

\noindent 她们来自内地，穷人家的姑娘，需要工作。走进收容院时不知道这是什么地方。一个星期内离开——或者留下，然后变得坚硬。

\noindent 其中一个，1964年的，不一样。

\noindent 她叫伊莎贝尔。19岁。黑头发，大眼睛，递食物时手在发抖。

\noindent 第一个星期，她经过我的牢房门口，看见了日记。

\noindent \textit{“您写什么？”}

\noindent \textit{“真相。”}

\noindent 她没有听懂。但她继续经过。

\noindent 每天下午，她送来食物。其他护工把盘子从门上的洞口扔进来，她却打开门，递到我手里。

\noindent \textit{“您不危险，”有一天她说。}

\noindent \textit{“你怎么知道？”}

\noindent \textit{“危险的人不写作。危险的人不思考。您写作，您思考。您只是悲伤。”}

\noindent \textit{“悲伤？”}

\noindent 她垂下眼睛。

\noindent \textit{“您有一双我见过的最悲伤的眼睛。”}

\noindent \textit{“它是黄色的。”}

\noindent 她摇了摇头。

\noindent \textit{“黄色是颜色。悲伤是颜色背后的东西。”}

\noindent 她走了。

\noindent 那一夜，怪物开口了：

\noindent \textit{“她喜欢你。”}

\noindent \textit{“闭嘴。”}

\noindent \textit{“她想救你。”}

\noindent \textit{“在这里谁也救不了谁。”}

\noindent \textit{“她想试一试。”}

\noindent \textit{“……她会走的。”}

\noindent \textit{“你怎么知道？”}

\noindent \textit{“所有人都会走。每一个进来的人，都会离开。要么走着出去。要么躺着出去。”}

\vspace{0.5em}\noindent\rule{\textwidth}{0.5pt}\vspace{0.5em}

\noindent 伊莎贝尔待了三个月。

\noindent 有一天，她不再来了。

\noindent 送饭的护工说：\textit{“辞职了。说她看不下去。”}

\noindent 我问：\textit{“看不下什么？”}

\noindent 他看着我。

\noindent \textit{“看你们。”}

\noindent 伊莎贝尔走后，我合上日记，一个星期没有打开。

\noindent 不是因为缺少词语。
是因为缺少希望。

\noindent 如果连想要帮忙的人都走了，那就什么都不值得了。

\noindent 收容院不是医院。
是一个坑。
而所有掉进坑里的东西——都被遗忘。

\vspace{0.5em}
\noindent\textit{\textbf{↪ 路线:} \rota{MEM-25} \rotasep \rota{MEM-12}}
